\scanpage{114}
bijection canonique
\[
\Gamma_{i\rho}\cap\Sigma_{\mathbb R}\simeq I\setminus\{i\}
\]
en associant à $\forall x\in\Gamma_{i\rho}\cap\unsure{E_0|\Sigma_{\mathbb R}}$ l'élément de $I\setminus\{i\}$ le plus proche. Ces deux pts sont \unsure{diamétralmt|diamétralement} opposés sur $\Gamma_{i\rho}$, donc on trouve un isomorphisme canonique
\[
\Gamma_{i\rho}\simeq\deleted{\mathbb Z}\cdot\mathbb U\wedge_{\{\pm1\}}(I\setminus\{i\})
\]
ie. $\Gamma_{i\rho}$ est déduit du $\{\pm1\}$-torseur $I-\{i\}$ par l'inclusion
\[
\deleted{(\mathbb Z\times\unsure{\mu_2|\{\pm1\}})}\inserted{\{\pm1\}}\lhook\joinrel\longrightarrow\mathbb U.
\]

\deleted{la donnée de} La courbe $X_G$ du type $G$ s'obtient en prenant la normalisée
\[
\widetilde X_G=\coprod_{s\in\unsure{G|S}}\Sigma_{\widehat{\vec A}_s}
\]
par recollement à partir de l'involution $\sigma$ sur $\widehat{\vec A}=\coprod_{s\in S}\widehat{\vec A}_s$. On a pour $\forall x\in\widehat A$, au dessus de $s\in S$
\[
T_x\simeq\mathbb U\wedge_{\{\pm1\}}\{\widehat{\vec A}_s\setminus x\}
\qquad\text{si }x\in I\simeq\widehat{\vec A}^{\sigma}
\]
\inserted{(en identifiant les éléments de $I$ à leurs images dans $\widehat A$)}
\[
T_x\simeq\mathbb U\wedge_{\{\pm1\}}\bigwedge_{\widetilde x\in\widehat{\vec A}_x}
\bigl(\widehat{\vec A}_{\unsure{\sigma|o}(\widetilde x)}\setminus\{\widetilde x\}\bigr)
\qquad\text{si }x\in A
\]
\figure{\begin{tikzpicture}[x=1.05cm,y=1.05cm]
\coordinate (a) at (-2,-.65);\coordinate (b) at (-1.45,.12);\coordinate (c) at (-.65,.17);\coordinate (d) at (.55,.12);\coordinate (e) at (1.12,.80);\coordinate (f) at (1.48,.08);\coordinate (g) at (.85,-.50);\coordinate (h) at (1.22,1.40);
\draw (a)--(b)--(c)--(a);\draw (c)--(d)--(e)--(f)--(g)--(d);\draw (e)--(h);
\draw (h) .. controls (.75,1.85) and (1.00,2.22) .. (1.40,2.02) .. controls (1.90,1.85) and (1.57,1.56) .. (h);
\draw[->] (a)--(-2.12,.20);\draw[->] (b)--(-1.60,1.38);\draw[->] (f)--(2.18,.08);\draw[->] (g) .. controls (.91,-.88) and (.83,-1.03) .. (.80,-1.20);
\fill (a) circle (.024);\fill (b) circle (.030);\fill (c) circle (.023);\fill (d) circle (.030);\fill (h) circle (.035);
\end{tikzpicture}}
\uline{NB} Si $x$ est une boucle, on trouve $T_x\simeq\mathbb U$ iso. can.
\mspage{113}
